\documentclass[11pt,letterpaper]{article}
\usepackage[margin=1in]{geometry}
\usepackage[T1]{fontenc}
\usepackage{lmodern}
\usepackage{amsmath,amssymb,amsthm}
\usepackage{booktabs,array}
\usepackage{xcolor}
\usepackage{hyperref}
\hypersetup{colorlinks=true,linkcolor=blue!45!black,citecolor=blue!45!black,
  urlcolor=blue!45!black,pdftitle={From Scalar Certificates to Spectral Bounds},
  pdfauthor={Carl P. Zimmerman}}
\urlstyle{same}
\setlength{\emergencystretch}{2em}
\newcommand{\paperdoi}{\href{https://doi.org/10.5281/zenodo.22895866}{doi:10.5281/zenodo.22895866}}
\newcommand{\R}{\mathbb R}
\newcommand{\C}{\mathbb C}
\newcommand{\Z}{\mathbb Z}
\newcommand{\SU}{\operatorname{SU}}
\newcommand{\Tr}{\operatorname{Tr}}
\newcommand{\Spec}{\operatorname{Spec}}
\newcommand{\essinf}{\operatorname*{ess\,inf}}
\newcommand{\gap}{\operatorname{gap}}
\newcommand{\dd}{\,\mathrm d}
\newtheorem{theorem}{Theorem}[section]
\newtheorem{proposition}[theorem]{Proposition}
\newtheorem{lemma}[theorem]{Lemma}
\newtheorem{corollary}[theorem]{Corollary}
\theoremstyle{definition}
\newtheorem{remark}[theorem]{Remark}
\numberwithin{equation}{section}

\title{From Scalar Certificates to Spectral Bounds:\\
Radial Pulsations, Discrete Hardy Forms,\\
and $\SU(N)$ Lattice Gauge Theory}
\author{Carl P. Zimmerman\\[2pt]\small Briar Creek Tech}
\date{September 22, 2026\\\small Version 1.0\quad---\quad\paperdoi}

\begin{document}
\maketitle

\begin{abstract}
A scalar inequality becomes a spectral statement only after its operator,
norm, domain, and comparison direction have been identified. We carry out
that identification for three families of certificates in a public research
repository. For adiabatic radial pulsations, an exact Newtonian work identity
and a relativistic positive-function comparison distinguish a homologous
trial quotient from the fundamental eigenvalue. The first post-Newtonian
correction is a positive structure integral; for an $n=3$ polytrope its
coefficient in the repository normalization is approximately $3.37342294$.
For a specified finite Dirichlet difference form, we obtain the complete
sine spectrum and the sharp lower bound uniform in box size, including the
shift caused by constant confinement. Direct variation of the proposed
scalar source action reveals an obstruction to identifying that form with
the physical fluctuation operator. For a normalized $\SU(N)$ lattice
Hamiltonian, min--max gives an untruncated single-plaquette gap bound.
An application of Yarotsky's weak-interaction theorem then gives a gap
at least $3g^2/16$, uniformly in $N$ and open volume, at sufficiently large
$g^2$ for each fixed spatial dimension. We state precisely which steps are
analytic, externally supplied, numerically checked, or Lean-certified.
No stellar mass cutoff, unique physical coupling, or continuum Yang--Mills
gap is inferred.
\end{abstract}

\section{Introduction and scope}

An algebraic expression can be positive without being an eigenvalue or a
lower bound on one. A positive Rayleigh quotient, for example, need not
exclude a negative eigenvalue. Similarly, a finite matrix calculation does
not by itself bound the spectrum of an untruncated operator, and a sharp
inequality for a chosen quadratic form does not identify that form with a
physical second variation. These distinctions motivate the present work.

The starting objects are three repository certificates labelled I13, I14,
and I15. Their scalar statements concern, respectively, pressure and
relativistic work, a Dirichlet difference energy, and an $N$-dependent
strong-coupling expression. The labels are provenance identifiers, not
additional assumptions. The purpose here is to give the operator statements
that can be supported, and to locate the physical identifications that
remain missing. The accompanying derivations and executable records are
collected in the source archive described in Section~\ref{sec:reproduce}.

The methods are established ones: radial variational principles,
completion of squares, finite Dirichlet diagonalization, compact-group
representation theory, and stability of gapped product systems. We claim
no novelty for those methods or for the existence of strong-coupling lattice
gaps. The specific contribution is an explicit normalization and domain
analysis connecting them to the stated certificates, together with
counterarguments to stronger identifications.

For radial pulsations, the differential equation and density convention
are tied to the versioned source of Saio et al.~\cite{saio}. For lattice
gauge theory, the character description is consistent with the
Hamiltonian treatment of Ligterink, Walet, and Bishop~\cite{ligterink};
the volume-uniform result uses Yarotsky's theorem~\cite{yarotsky} as an
external input. The discrete form and the scalar source-action comparison
are specified directly below.

\begin{table}[htbp]
\centering\small
\begin{tabular}{@{}p{0.13\linewidth}p{0.40\linewidth}p{0.40\linewidth}@{}}
\toprule
Certificate & Supported spectral statement & Additional input still required\\
\midrule
I13 & Upper trial and lower positive-function comparisons for a specified
radial operator & Stellar profiles, equation of state, relevant inner
boundary, and controlled relativistic evaluation for a mass cutoff\\[4pt]
I14 & Complete constant-confinement finite spectrum and sharp bound uniform
in box size & Physical action-to-form map and an independent marginality
condition for selecting a coupling\\[4pt]
I15 & Single-plaquette gap for $g^2\ge2$; open-volume and canonical GNS gap
for $g^2\ge X_d$ & Numerical $X_d$, any extension to the claimed smaller
window, and a separate continuum analysis\\
\bottomrule
\end{tabular}
\caption{Scope of the three operator comparisons. The constant $X_d$ is
finite but is not numerically evaluated in this work.}
\label{tab:scope}
\end{table}

We use different symbols where the original certificates reused a
parameter: $\chi=2GM/(Rc^2)$ denotes stellar compactness, $L$ the number of
edges in a one-dimensional Dirichlet box, and $x=g^2$ the lattice gauge
coupling. All finite-dimensional eigenvalues are relative to the explicitly
specified norm. In an infinite-dimensional problem, a statement about the
bottom of the spectrum does not imply that the bottom is attained unless
this is separately assumed or proved.

\section{Radial pulsations: trial values and spectral comparisons}
\label{sec:radial}

\subsection{Equilibrium and domain conventions}

Let $P(r)$, $\rho(r)$, $m(r)$, and $\Gamma(r)$ denote pressure, density,
enclosed mass, and adiabatic exponent in a spherical equilibrium of radius
$R$. Write a radial displacement as $\xi(r)e^{i\Omega t}$ and put
$\eta=\xi/r$. Assume $P,\rho,\Gamma>0$ in the interior, sufficient local
regularity for the equations below, a regular center, and a vacuum surface
with the natural pressure boundary condition. Integrations by parts are
used only for functions for which the displayed endpoint terms vanish.
One may first work on admissible smooth test functions and then take the
closure in the radial form norm. Whenever the homologous displacement
$\eta=1$ is used, its membership in that domain is required.

We assume the resulting form has a lower-bounded self-adjoint realization.
The notation $\lambda_0$ initially means its spectral infimum; statements
requiring a lowest eigenfunction additionally assume attainment. These
conditions matter at singular endpoints. In particular, a star containing
a central black hole has an inner-boundary problem that is not covered by
the regular-center hypothesis. Rotation, nonadiabatic driving, accretion,
and dissipative boundary terms are also outside this model.

\subsection{The exact Newtonian work identity}

Set $K=3\Gamma-4$. The Newtonian radial equation in fractional displacement
is
\begin{equation}
 (\Gamma P r^4\eta')'+r^3(KP)'\eta
       +\Omega^2\rho r^4\eta=0.
 \label{eq:newton-ode}
\end{equation}
Define
\begin{align}
 J[\eta]&=\int_0^R\rho r^4\eta^2\dd r,\label{eq:inertia}\\
 \mathcal N_N[\eta]&=\int_0^R
       \bigl(\Gamma P r^4\eta'^2-r^3(KP)'\eta^2\bigr)\dd r.
 \label{eq:newton-work}
\end{align}
Multiplying~\eqref{eq:newton-ode} by $\eta$ shows that an eigenfunction has
$\Omega^2=\mathcal N_N[\eta]/J[\eta]$. For a general admissible nonzero
$\eta$, this ratio is a trial quotient.

\begin{proposition}[Newtonian square identity]
Suppose the boundary terms in~\eqref{eq:newton-work} and
$[r^3KP\eta^2]_0^R$ vanish. With $\delta=\Gamma-4/3$,
\begin{equation}
 \mathcal N_N[\eta]
 =\int_0^R P\left\{\frac43r^4\eta'^2
       +\delta r^2(r\eta'+3\eta)^2\right\}\dd r.
 \label{eq:newton-square}
\end{equation}
This identity permits spatially varying $\Gamma$.
\end{proposition}
\begin{proof}
Integration of the differentiated pressure factor gives
\[
 \mathcal N_N[\eta]=\int_0^R P
  \bigl(\Gamma r^4\eta'^2+2Kr^3\eta\eta'+3Kr^2\eta^2\bigr)\dd r.
\]
Insert $K=3\delta$ and $\Gamma=4/3+\delta$, and complete the square.
No derivative of $\Gamma$ is discarded in this operation.
\end{proof}

Consequently, $\Gamma\ge4/3$ pointwise implies nonnegativity of this
Newtonian form. If $\Gamma>4/3$ on a set of positive measure, a zero-energy
function must have $\eta'=0$ and then $\eta=0$. When the spectral bottom
is attained, this rules out a zero lowest eigenvalue and gives
$\lambda_0>0$. Without attainment, positivity on each nonzero vector alone
would not imply a positive uniform lower bound. If $\Gamma=4/3$
identically, admissible homology is an exact zero mode.

Let
\begin{equation}
 W_\beta=\int_0^R KP r^2\dd r,\qquad
 W_{\rm gr}=\int_0^R P r^2\dd r,\qquad
 J=\int_0^R\rho r^4\dd r.
 \label{eq:moments}
\end{equation}
The homologous trial has
\begin{equation}
 Q_N[1]=\frac{3W_\beta}{J},\qquad
 \lambda_0\le\frac{3W_\beta}{J}.
 \label{eq:homology}
\end{equation}
The I13 scalar expression has Newtonian part $W_\beta/I$. Its equality
with this trial quotient therefore requires $I=J/3$. If instead an angular
factor is included in a conventional modal inertia, the numerator must
carry the corresponding factor as well. The factor three in
\eqref{eq:homology} cannot be removed by terminology.

There is also a direct test for when homology is an eigenmode. Acting on
$1$ with the differential operator in~\eqref{eq:newton-work} gives
$-r^3(KP)'$. Thus the quotient can be an eigenvalue of that constant
function only if
\begin{equation}
 -\frac{(KP)'}{\rho r}\quad\hbox{is constant almost everywhere.}
 \label{eq:homology-residual}
\end{equation}
For constant $\Gamma\ne4/3$, hydrostatic balance
$P'=-Gm\rho/r^2$ turns this into constancy of $m/r^3$. A nonuniform-density
polytrope therefore generally violates the required condition. This is an
exact obstruction to a universal identification of the scalar trial
formula with the fundamental eigenvalue.

For constant $K>0$, the same equation supplies a genuine enclosure:
\begin{equation}
 K\essinf_{0<r<R}\frac{Gm(r)}{r^3}
 \ \le\lambda_0\le\frac{3K W_{\rm gr}}{J}.
 \label{eq:newton-enclosure}
\end{equation}
Indeed, the gradient part of~\eqref{eq:newton-work} is nonnegative and the
remaining part is the weighted integral of $KGm/r^3$. For decreasing
density, the mean interior density decreases with $r$, giving the lower
endpoint $KGM/R^3$.

\subsection{The relativistic operator and a lower comparison}

Adopt the metric convention
\begin{equation}
 ds^2=-e^{2a(r)}c^2dt^2+e^{2b(r)}dr^2+r^2d\Omega_{\mathbb S^2}^2,
 \qquad e^{-2b}=1-\frac{2Gm}{rc^2}.
\end{equation}
The density convention is the one in Saio et al.~\cite{saio}, equations
(2)--(5), without an additional internal-energy correction. This is a
specified operator convention, not an assertion that it supplies a complete
equation of state for every stellar application. We require $2Gm/(rc^2)<1$
throughout the equilibrium. The equilibrium identities used below are
\begin{equation}
 P'=-(P+\rho c^2)a',\qquad
 a'+b'=\frac{4\pi G r e^{2b}(P+\rho c^2)}{c^4}.
 \label{eq:tov-identities}
\end{equation}

Set $z=e^{-a}r^2\xi=e^{-a}r^3\eta$ and $\Omega=c\sigma$. Rearrangement
of the radial equation gives
\begin{equation}
 -(\Pi z')'+Vz=\sigma^2Wz,
 \label{eq:gr-sl}
\end{equation}
where
\begin{align}
 \Pi&=\frac{e^{3a+b}\Gamma P}{r^2},
 &W&=\frac{e^{a+3b}(P+\rho c^2)}{r^2},\label{eq:gr-coeff}\\
 V&=\frac{e^{3a+b}}{r^2}
 \left(\frac{4P'}r+
 \frac{8\pi G e^{2b}P(P+\rho c^2)}{c^4}
 -\frac{P'^2}{P+\rho c^2}\right).
 \label{eq:gr-potential}
\end{align}
In particular, the dimensional variational quotient is
\begin{equation}
 Q_{\rm GR}[z]
 =c^2\frac{\int_0^R(\Pi z'^2+Vz^2)\dd r}
               {\int_0^R Wz^2\dd r},\qquad
 \lambda_0\le Q_{\rm GR}[z].
 \label{eq:gr-rayleigh}
\end{equation}

\begin{proposition}[Positive-function comparison]
Let $u>0$ in the interior, with enough regularity for
$-(\Pi u')'+Vu$ to be defined, and suppose
$[\Pi u'z^2/u]_0^R=0$ for the test functions considered. Put
\[
 q_u=\frac{c^2[-(\Pi u')'+Vu]}{Wu}.
\]
Then $\lambda_0\ge\essinf q_u$. If $u$ is also an admissible trial
function, then
\begin{equation}
 \essinf q_u\le\lambda_0\le Q_{\rm GR}[u].
 \label{eq:gr-enclosure}
\end{equation}
\end{proposition}
\begin{proof}
Write $z=u v$, expand its derivative, and integrate the cross term. This
gives the exact identity
\begin{equation}
 \int_0^R(\Pi z'^2+Vz^2)\dd r
 -\int_0^R\frac{-(\Pi u')'+Vu}{u}z^2\dd r
 =\int_0^R\Pi u^2\bigl((z/u)'\bigr)^2\dd r.
 \label{eq:ground-transform}
\end{equation}
The right side is nonnegative. Divide by $\int Wz^2$ and take the
variational infimum. The upper inequality is~\eqref{eq:gr-rayleigh}.
\end{proof}

Thus a negative trial value proves a negative spectral bottom, and a
positive lower endpoint proves positivity. A positive upper endpoint
alone proves neither. If $q_u$ is constant, the two endpoints coincide.
Endpoint compatibility in~\eqref{eq:ground-transform} is a hypothesis of
this conclusion, not a formal consequence of choosing $u>0$.

\subsection{A form suitable for computation and expansion}

Substituting $z=e^{-a}r^3\eta$ in~\eqref{eq:gr-rayleigh}, using
\eqref{eq:tov-identities}, and integrating the derivative
$4(e^{a+b}Pr^3\eta^2)'$ yields $Q_{\rm GR}=\mathcal N_{\rm GR}/\mathcal D_{\rm GR}$,
with
\begin{align}
 \mathcal N_{\rm GR}[\eta]
 ={}&\int_0^R e^{a+b}\Bigl\{\Gamma P r^4\eta'^2
       +2Pr^3(K-\Gamma r a')\eta\eta'\notag\\
 &\quad+\bigl[3KPr^2-(6\Gamma+2)Pr^3a'-2Pr^3b'
       +((\Gamma-1)P-\rho c^2)r^4a'^2\bigr]\eta^2\Bigr\}\dd r,
 \label{eq:stable-gr-form}\\
 \mathcal D_{\rm GR}[\eta]
 ={}&\int_0^R e^{-a+3b}(\rho+P/c^2)r^4\eta^2\dd r.
 \label{eq:stable-gr-mass}
\end{align}
The removed endpoint term vanishes under the regular-center and vacuum
surface conditions. The derivative form~\eqref{eq:gr-rayleigh} and the
expanded form~\eqref{eq:stable-gr-form} provide separate assembly routes
for a numerical implementation.

Consider a weak-field sequence about a Newtonian equilibrium with
$\Gamma-4/3=O(\varepsilon)$ and $Gm/(rc^2)=O(\varepsilon)$. At first
order, evaluate the leading integrals on the Newtonian background. For
homology, the result is
\begin{align}
 \mathcal N_{\rm GR}[1]&=3W_\beta-\frac{\mathcal H}{c^2}
                       +O(\varepsilon^2),\label{eq:pn-work}\\
 \mathcal H&=\int_0^R
   \left(8GPmr+8\pi GP\rho r^4+G^2\rho m^2\right)\dd r>0.
 \label{eq:structure-integral}
\end{align}
The order term in~\eqref{eq:pn-work} is measured in the pressure-work
scale; it is not a supplied numerical error bound. To see the coefficients,
use
\[
 a'=\frac{Gm}{r^2c^2}+O(\varepsilon^2/r),\qquad
 b'=\frac{G(4\pi r\rho-m/r^2)}{c^2}+O(\varepsilon^2/r)
\]
in~\eqref{eq:stable-gr-form}. At $\Gamma=4/3$, the terms linear in
$a',b'$ combine as $-8GPmr/c^2-8\pi GP\rho r^4/c^2$. The factor
$-\rho c^2r^4a'^2$ supplies the remaining term
$-G^2\rho m^2/c^2$. The other products enter at the next order under the
stated scaling. Since the numerator starts at first order and
$\mathcal D_{\rm GR}[1]=J+O(\varepsilon)$ in the inertia scale, its
denominator correction affects the quotient only at second order.

Define the profile-dependent coefficient
\begin{equation}
 C_{\rm gr}=\frac{\mathcal H R}{6GM W_{\rm gr}},\qquad I=J/3,
 \qquad \chi=\frac{2GM}{Rc^2}.
 \label{eq:cgr}
\end{equation}
Then the first-order homologous quotient takes exactly the scalar form
\begin{equation}
 Q_{\rm PN}[1]=\frac{W_\beta-C_{\rm gr}\chi W_{\rm gr}}{I}.
 \label{eq:pn-scalar}
\end{equation}
This identification concerns a trial quotient and a controlled-order
expansion. It is not exact linear dependence of a physical eigenvalue on
compactness. A varying equilibrium sequence also changes its moments and
its structure coefficient.

If an explicit bound
$|\mathcal N_{\rm GR}[1]-(3W_\beta-\mathcal H/c^2)|\le E$ is available,
then $3W_\beta-\mathcal H/c^2+E<0$ suffices for full-model instability.
Otherwise one must evaluate the exact numerator to make such a conclusion
at finite compactness. An asymptotic order symbol alone is insufficient
near a sign change.

For a uniform-density Newtonian star with $G=\rho=R=1$,
$m=4\pi r^3/3$ and $P=2\pi(1-r^2)/3$. Direct integration gives
$W_{\rm gr}=4\pi/45$ and $\mathcal H=304\pi^2/315$, hence
\begin{equation}
 \frac{\mathcal H R}{18GM W_{\rm gr}}=\frac{19}{42},
 \qquad C_{\rm gr}=\frac{19}{14}.
 \label{eq:uniform-cgr}
\end{equation}
For the generated $n=3$ polytrope, the corresponding numerical values are
\begin{equation}
 \frac{\mathcal H R}{18GM W_{\rm gr}}\simeq1.124474311979,
 \qquad C_{\rm gr}\simeq3.373422935937.
 \label{eq:n3-cgr}
\end{equation}
The first value agrees, after rounding, with the $1.1245$ benchmark in
Saio et al.~\cite{saio}. The factor-three distinction follows from
$K=3\Gamma-4$. In this normalization the previously used upper endpoint
$3.35$ is too low for that polytrope, and no profile-independent coefficient
band is justified by these two examples.

\section{A complete discrete Dirichlet spectrum}
\label{sec:discrete}

Let $L\ge2$, let $u=(u_1,\ldots,u_{L-1})\in\R^{L-1}$, and set
$u_0=u_L=0$. Define
\begin{align}
 U[u]&=\sum_{n=0}^{L-1}u_n^2,&
 D[u]&=\sum_{n=0}^{L-1}(u_{n+1}-u_n)^2,\notag\\
 A[u]&=\sum_{n=0}^{L-1}(u_{n+1}-\tfrac32u_n)^2,&
 E[u]&=A[u]-\alpha U[u]+\beta\sum_{n=0}^{L-1}q_nu_n^2.
 \label{eq:discrete-form}
\end{align}
Here $\alpha,\beta$ are real parameters. The interpretation
$\alpha=\kappa^2$, $\beta=\mu^2$ is an additional restriction. We study
the matrix $H$ defined by $E[u]=u^THu$ relative to the norm $U$.
The ordinary coordinate Hessian of $E$ is $2H$; this convention has no
effect on signs but does affect quoted eigenvalues.

\begin{theorem}[Constant-confinement spectrum]
If $q_n=q_0$ and $b_0=1/4-\alpha+\beta q_0$, the eigenvalues of $H$ are
simple and equal to
\begin{equation}
 \lambda_j=6\sin^2\left(\frac{j\pi}{2L}\right)+b_0,
 \qquad 1\le j<L.
 \label{eq:sine-spectrum}
\end{equation}
The corresponding eigenvectors are $v_j(n)=\sin(j\pi n/L)$. In
particular,
\begin{equation}
 \min_{u\ne0}\frac{E[u]}{U[u]}
 =6\sin^2\left(\frac{\pi}{2L}\right)+b_0.
 \label{eq:finite-min}
\end{equation}
\end{theorem}
\begin{proof}
Telescoping the endpoint values gives
\[
 2\sum_{n=0}^{L-1}(u_{n+1}-u_n)u_n=-D[u].
\]
Expanding $u_{n+1}-\tfrac32u_n=(u_{n+1}-u_n)-\tfrac12u_n$
therefore yields
\begin{equation}
 A[u]=\frac32D[u]+\frac14U[u],\qquad
 E[u]=\frac32D[u]+b_0U[u].
 \label{eq:hardy-identity}
\end{equation}
The matrix of $D$ has diagonal $2$ and nearest off-diagonal entries $-1$.
For $\theta_j=j\pi/L$, the addition formula gives
\[
 2v_j(n)-v_j(n-1)-v_j(n+1)
 =\bigl(2-2\cos\theta_j\bigr)v_j(n).
\]
The boundary values vanish, and $v_j(1)>0$ for $1\le j<L$, so these
vectors are nonzero. Their eigenvalues increase strictly because
$0<j\pi/(2L)<\pi/2$. Symmetry of the matrix makes eigenvectors at distinct
eigenvalues orthogonal. There are $L-1$ of them in a space of dimension
$L-1$, proving completeness. Expanding any $u$ in this basis makes its
Rayleigh quotient a weighted average of the eigenvalues, proving
\eqref{eq:finite-min}.
\end{proof}

Thus the form on one fixed box is nonnegative precisely when
\begin{equation}
 \alpha\le\frac14+\beta q_0+
             6\sin^2\left(\frac{\pi}{2L}\right).
 \label{eq:finite-wall}
\end{equation}
Strict inequality gives positive definiteness; equality gives the
one-dimensional kernel spanned by $v_1$. If the inequality is reversed,
amplitude scaling makes the unrestricted quadratic energy unbounded
below, although its unit-norm Rayleigh quotient still has a finite minimum.
The boxes $L=0,1$ have no nonzero interior modes and are excluded from the
eigenvalue assertion.

\begin{theorem}[Sharp bound uniform in box size]
For constant confinement,
\begin{equation}
 \inf_{L\ge2}\ \inf_{u\ne0}\frac{E[u]}{U[u]}=b_0.
 \label{eq:uniform-floor}
\end{equation}
No nonzero finite-box vector attains this infimum. A lower bound
$E\ge\gamma U$ holds on every finite box if and only if $\gamma\le b_0$.
\end{theorem}
\begin{proof}
Equation~\eqref{eq:hardy-identity} gives $E\ge b_0U$. Equality on a
nonzero vector would require $D=0$, hence a constant vector with zero
endpoints, a contradiction. For sharpness take $u_n=1$ at all interior
sites. Then $U=L-1$, $D=2$, and
\begin{equation}
 \frac{E[u]}{U[u]}=b_0+\frac3{L-1}.
 \label{eq:slab}
\end{equation}
For every $\epsilon>0$, an integer $L>1+3/\epsilon$ gives a quotient
below $b_0+\epsilon$. No larger uniform bound can hold.
\end{proof}

There are three distinct regimes. If $b_0>0$, a positive lower bound is
uniform in all boxes. If $b_0=0$, every fixed finite box is positive but
there is no positive bound uniform in box size. If $b_0<0$, sufficiently
large boxes have negative modes: the explicit slab is negative whenever
$L>1+3/(-b_0)$. Slabs establish sharpness; they are generally not the
lowest sine eigenvectors on a fixed box.

In the square-parameter convention without confinement, all-box
nonnegativity is $|\kappa|\le1/2$. Restricting $\kappa\ge0$ changes this
to $0\le\kappa\le1/2$, not to the single selected value $1/2$.
With constant confinement it becomes
$\kappa^2\le1/4+\mu^2q_0$. For example, taking
$\kappa^2=\mu^2=q_0=1$ gives
$E=(3/2)D+U/4\ge U/4$ despite $\kappa=1$. Positive confinement therefore
rules out a confinement-independent wall at $\kappa=1/2$.

For variable $q_n\ge q_0$ and $\beta\ge0$, the same identity gives
$E\ge b_0U$. This is optimal uniformly over the entire class
$q\ge q_0$, because the class contains the constant profile. It need
not be sharp for a specified nonconstant profile, and a sign assumption
on $\beta$ is essential for this comparison.

\section{Why the scalar source action does not identify the discrete form}
\label{sec:source-bridge}

The exact discrete theorem does not establish that its energy is a
physical fluctuation energy. We can test the proposed identification
directly against the repository kinetic action, rather than assume that
the repeated occurrence of $1/4$ supplies it. The source paths and their
pinned hashes are recorded with the I14 computations~\cite{archive}.

On a fixed flat spatial background consider the radial functional
\begin{equation}
 \mathcal F[\phi]=4\pi\int r^2\Lambda^4f(K)\dd r,
 \qquad K=\frac{\phi_r^2}{2\Lambda^4},
 \label{eq:scalar-action}
\end{equation}
with $\Lambda>0$ a scale. The G155 source specifies, for $y=\sqrt K>0$,
\begin{equation}
 f(K)=y^2-2\log(1+y)-\frac2{1+y}+1,
 \qquad f_K=\mu(y)=\frac{y(y+2)}{(1+y)^2}.
 \label{eq:interpolant}
\end{equation}
The symbol $\mu(y)$ here is the source's interpolation function; it is
unrelated to the confinement parameter in Section~\ref{sec:discrete}.

For $\phi=\phi_0+\epsilon h$, direct Taylor expansion shows that the
coefficient of $\epsilon^2$ is
\begin{equation}
 \mathcal Q[h]=2\pi\int r^2a_r(r)h_r^2\dd r,
 \qquad a_r=f_K+2Kf_{KK}.
 \label{eq:source-second}
\end{equation}
Indeed, the first and second perturbations of $K$ are
$\phi_{0,r}h_r/\Lambda^4$ and $h_r^2/(2\Lambda^4)$, respectively.
Substitution in $f$ gives~\eqref{eq:source-second}; the full second
variation is twice $\mathcal Q$. In this fixed-background calculation,
shift symmetry produces a derivative-only quadratic form in $h$.

For~\eqref{eq:interpolant},
\begin{equation}
 a_r(y)=\mu(y)+y\mu'(y)
       =\frac{y(y^2+3y+4)}{(1+y)^3}>0.
 \label{eq:radial-coeff}
\end{equation}
Take the proposed logarithmic profile $\phi_0=C\log r$, with $C>0$, and
write $\ell=C/(\sqrt2\Lambda^2)$, so $y=\ell/r$. With $r=e^t$,
\begin{equation}
 \frac{\mathcal Q[h]}{2\pi}=\int w(t)h_t^2\dd t,
 \qquad w(t)=r a_r(\ell/r)
       =\ell\frac{y^2+3y+4}{(1+y)^3}.
 \label{eq:log-weight}
\end{equation}
In the deep regime $y\to0$, this weight tends to the constant $4\ell$.
It does not tend to the canonical radial weight proportional to $e^t$.

For any positive weight, the unit-weight variable $v=\sqrt w\,h$ gives
\begin{equation}
 \frac{\mathcal Q[h]}{2\pi}=\int(v_t-sv)^2\dd t
 =\int\bigl(v_t^2+(s^2+s')v^2\bigr)\dd t,
 \qquad s=\tfrac12(\log w)',
 \label{eq:weight-potential}
\end{equation}
under compact support or vanishing endpoint terms. A constant radial
$a_r$ has $w\propto e^t$, $s=1/2$, and potential $1/4$; by contrast,
$a_r\propto r^{-1}$ has constant $w$ and potential zero. More generally,
$a_r\propto r^{-p}$ gives potential $(1-p)^2/4$. The weight, rather
than the mere appearance of a logarithmic profile, determines this number.

The exact interpolant gives
\begin{align}
 s(y)&=\frac{y(y^2+4y+9)}{2(1+y)(y^2+3y+4)},\notag\\
 V_{\rm dress}(y)&=s(y)^2-y\frac{ds}{dy}\notag\\
 &=\frac{y(y^5+8y^4+42y^3+64y^2+17y-72)}
              {4(1+y)^2(y^2+3y+4)^2}.
 \label{eq:dressed-source}
\end{align}
It tends to zero as $y\to0$ and to $1/4$ as $y\to\infty$. It is not a
constant $1/4-\kappa^2$. A locally negative value of this dressed
potential also does not imply a negative mode, since the original form
is an integral of positive weighted squares.

There is a second obstruction: the proposed logarithmic profile is not
an exact source-free background for the full interpolant. Its radial flux
is
\begin{equation}
 r^2\mu(y)\phi_{0,r}=C\ell\frac{y+2}{(1+y)^2},\qquad
 r\frac{d}{dr}\bigl(r^2\mu(y)\phi_{0,r}\bigr)
 =C\ell\frac{y(y+3)}{(1+y)^3}>0.
 \label{eq:flux-obstruction}
\end{equation}
Source-free equilibrium requires constant flux. The leading approximation
$\mu(y)=2y$ does give constant flux $2C\ell$, so the profile is compatible
with that approximate equation, but the exact and approximate backgrounds
must be distinguished.

Finally, another cited repository source uses $f(K)=K-1/(1+K)$ with a
different invariant convention. Its derivative at zero is $2$, whereas
the derivative in~\eqref{eq:interpolant} tends to zero. A sign convention
alone cannot identify these two functions. The related isothermal-fluid
calculations also involve a different displacement, norm, and endpoint
problem; local solutions of an ordinary differential equation do not
determine the spectrum of a fixed self-adjoint domain.

These calculations concern a spatial functional on a fixed background.
They do not decide the time kinetic sign or the stability of a fully
coupled matter--gravity theory. What they do establish is the failure of
the proposed direct map to~\eqref{eq:discrete-form}. A physical application
requires a specified action, its actual background, the constrained
quadratic dynamics, its Hilbert norm and endpoints, and a mesh and parameter
map. Selection of a unique coupling further requires a physical saturation
condition beyond nonnegativity.

\section{Spectral bounds for the \texorpdfstring{$\SU(N)$}{SU(N)} lattice Hamiltonian}
\label{sec:gauge}

\subsection{Hamiltonian, normalization, and Casimir minimum}

Fix $d\ge2$, $N\ge2$, $x=g^2>0$, and $0\le b_N\le2N$, at lattice
spacing one. Each positive lattice link has Hilbert space
$L^2(\SU(N))$ with normalized Haar measure. Reverse links are inverses,
not independent factors. Let $C_l$ be the positive Casimir with generator
normalization $\Tr(t_at_b)=\delta_{ab}/2$, and put
\begin{equation}
 H_N(x)=\frac x2\sum_l C_l+\frac{b_N}{x}\sum_p(1-T_p),
 \qquad T_p=\frac{\operatorname{Re}\Tr U_p}{N}.
 \label{eq:lattice-H}
\end{equation}
The trace is a real multiplication operator with $-1\le T_p\le1$.
The physical space consists of vectors invariant under all vertex gauge
transformations, with no external charges. The repository normalization
$b_N=2$, as well as $b_N=N$ and $b_N=2N$, is covered by the assumptions.
For finitely many links the electric Casimir sum has compact resolvent;
bounded magnetic multiplication preserves self-adjointness on its domain
and compact resolvent.

\begin{lemma}[Minimum nonzero Casimir]
The least nonzero Casimir eigenvalue in $L^2(\SU(N))$ is
\begin{equation}
 C_F=\frac{N^2-1}{2N}\ge\frac34.
 \label{eq:casimir-floor}
\end{equation}
\end{lemma}
\begin{proof}
With roots of squared length two, the fundamental weights have Gram matrix
\[
 (\omega_i,\omega_j)=\frac{\min(i,j)(N-\max(i,j))}{N}>0.
\]
This follows by writing
$\omega_i=\sum_{k\le i}e_k-(i/N)\sum_{k\le N}e_k$ in the hyperplane of
zero coordinate sum. An irreducible representation has highest weight
$\lambda=\sum_i a_i\omega_i$, with $a_i$ nonnegative integers, and
$C_2(\lambda)=(\lambda,\lambda+2\varrho)/2$, where
$\varrho=\sum_i\omega_i$. If $\lambda\ne0$, choose $a_i\ge1$ and
write $\lambda=\omega_i+\nu$ with $\nu$ dominant. Then
\[
 C_2(\lambda)-C_2(\omega_i)
 =\tfrac12(\nu,\nu)+(\nu,\omega_i+\varrho)\ge0.
\]
Moreover,
\[
 C_2(\omega_i)=\frac{i(N-i)(N+1)}{2N}\ge C_F,
\]
because $i(N-i)-(N-1)=(i-1)(N-i-1)\ge0$. The fundamental representation
attains equality. Peter--Weyl decomposition supplies all these sectors,
so no representation cutoff is involved.
\end{proof}

\subsection{An untruncated single-plaquette bound}

On an open four-edge square, gauge fixing a spanning tree leaves one
holonomy modulo conjugation. The physical Hilbert space is therefore the
space of class functions on $\SU(N)$, with irreducible characters as an
orthonormal basis. This standard reduction is also described in
\cite{ligterink}. Each of the four link Casimirs acts on a character
$\chi_R$ by $C_2(R)$, so
\begin{equation}
 H_E\chi_R=2xC_2(R)\chi_R,
 \qquad E_0(H_E)=0,\quad E_1(H_E)=2xC_F.
 \label{eq:one-loop-electric}
\end{equation}
Eigenvalues here and below are ordered with multiplicity.

\begin{theorem}[Single plaquette]
For the open square with Hamiltonian~\eqref{eq:lattice-H},
\begin{equation}
 E_1(H)-E_0(H)\ge2xC_F-\frac{b_N}{x}
 \ge x\frac{N^2-1}{N}-\frac{2N}{x}.
 \label{eq:one-plaquette-bound}
\end{equation}
In particular, for every $N\ge2$ and $x\ge2$ this gap is at least one.
\end{theorem}
\begin{proof}
The magnetic operator is nonnegative. The full min--max principle and
\eqref{eq:one-loop-electric} give $E_1(H)\ge2xC_F$. The normalized
constant function is an admissible trial vector. Haar integration of the
nontrivial fundamental character vanishes, so its expectation is
$\langle1,H1\rangle=b_N/x$, giving $E_0(H)\le b_N/x$. Subtract the two
inequalities. Finally, for $x\ge2$,
\[
 x\frac{N^2-1}{N}-\frac{2N}{x}
 \ge\frac{N^2-2}{N}=N-\frac2N\ge1.
\]
\end{proof}

The proof acts on the full physical character space. The constant electric
vacuum is used only as a trial vector: when $b_N>0$,
$H1=(b_N/x)(1-T)$ is nonconstant. In particular, subtracting its expectation
does not subtract the interacting ground energy. At $x=2$, the bound using
$b_N=2$ is $2N-2/N-1$, giving $2$ for $\SU(2)$ and $13/3$ for $\SU(3)$.
These values, and the weaker values $1$ and $7/3$ obtained from $b_N\le2N$,
are lower bounds, not exact gaps.

\subsection{A common strong-coupling window for all $N$ and open volumes}

We use the following external input from Yarotsky~\cite{yarotsky},
Theorems 1--3, pp.~2--4. At each site of $\Z^d$ allow an arbitrary,
possibly infinite-dimensional Hilbert space with a nonnegative
self-adjoint onsite operator, simple zero vacuum, and gap at least one.
Add bounded self-adjoint interactions $\phi_z$ supported in translates of
one fixed finite set $S$. There are positive $c_1(S),c_2(S)$ such that,
if $\epsilon_*:=\sup_z\|\phi_z\|<c_1$, finite empty-boundary restrictions
have a unique ground vector and gap at least $1-c_2\epsilon_*$. The
canonical thermodynamic state and its GNS Hamiltonian exist, with the same
gap estimate. The constants depend only on $S$, not on site dimension or
the upper onsite spectrum. The cited paper reviews the proof in Section~2;
its cluster expansion is an external theorem, not reproduced or
Lean-formalized here.

\begin{theorem}[Uniform open-lattice strong-coupling bound]
For each fixed $d\ge2$ there exists a finite $X_d$, independent of $N$,
$b_N$, and finite open volume, such that for $N\ge2$, $0\le b_N\le2N$,
and $x\ge X_d$, the following statements hold.
\begin{enumerate}
\item Every finite open subgraph of the cubic lattice, with selected
elementary plaquettes, has a unique physical ground vector and
\begin{equation}
 H-E_0\ge\frac{3x}{16}\quad\hbox{on its physical vacuum complement}.
 \label{eq:finite-many-gap}
\end{equation}
If that complement is zero, the spectral-exclusion statement is vacuous.
\item The canonical GNS Hamiltonian for the homogeneous infinite
interaction has a simple zero eigenvalue and
\begin{equation}
 \Spec(H_\infty)\subset\{0\}\cup[3x/16,\infty).
 \label{eq:infinite-gap}
\end{equation}
The same conclusions hold on the cyclic space generated by local
gauge-invariant observables.
\end{enumerate}
\label{thm:uniform-gap}
\end{theorem}
\begin{proof}
Group the $d$ positive outgoing links at $z$ into the site space
\[
 \mathcal H_z=L^2(\SU(N)^d),\qquad
 h_z=C_F^{-1}\sum_{i=1}^d C_{(z,i)}.
\]
The product of constant functions is its unique vacuum, and
\eqref{eq:casimir-floor} makes the onsite gap exactly one. Divide energies
by $s=xC_F/2$ and remove the scalar magnetic constant. With
$m_d=\binom d2$, group plaquettes at their lower anchor:
\begin{equation}
 \phi_z=-\frac{2b_N}{x^2C_F}\sum_{i<j}T_{(z,ij)}.
 \label{eq:site-interaction}
\end{equation}
The four link factors lie in the site blocks $z,z+e_i,z+e_j$, so all
these interactions have the same allowed support
$S=\{0,e_1,\ldots,e_d\}$. Since $\|T_p\|\le1$,
\begin{equation}
 \epsilon_*\le\frac{2b_Nm_d}{x^2C_F}
 \le\frac{A_d}{x^2},\qquad A_d=\frac{32}{3}\binom d2.
 \label{eq:small-parameter}
\end{equation}
Indeed, $b_N/C_F\le4N^2/(N^2-1)\le16/3$. Thus the small parameter is
uniform in $N$, despite the growing representation content of each site.

Choose
\begin{equation}
 X_d=\max\left\{1,\sqrt{\frac{2A_d}{c_1(S)}},
                       \sqrt{2c_2(S)A_d}\right\}.
 \label{eq:threshold}
\end{equation}
For $x\ge X_d$, we have $\epsilon_*\le c_1/2<c_1$ and
$c_2\epsilon_*\le1/2$. The external theorem gives a normalized gap at
least $1/2$. Restoring units yields $xC_F/4\ge3x/16$ and subtracts the
actual interacting ground energy. We next verify the boundary and gauge
statements needed to apply this estimate in the announced spaces.

For a finite open subgraph, retain $L^2(\SU(N))$ only on its actual
positive links and replace each absent factor by $\C$. An empty site has
$h_z=0$ on a one-dimensional space; the gap condition on its zero-dimensional
vacuum complement holds vacuously. Let~\eqref{eq:site-interaction} include
only selected elementary plaquettes. Pad the finite site set so that it
contains $z+S$ for every active anchor. The resulting inhomogeneous model
is exactly an allowed finite restriction with the same interaction range
and norm bound. Dummy factors introduce neither link degrees of freedom
nor boundary charges.

Every term of~\eqref{eq:lattice-H} commutes with the finite product of
vertex gauge groups, including all endpoints of present links. The unique
unconstrained ground vector spans a continuous one-dimensional
representation of that product. A finite product of $\SU(N)$ groups has
no nontrivial continuous characters, so this vector is gauge invariant.
The physical space reduces the Hamiltonian and contains its ground vector.
Restricting the excited-space quadratic-form inequality proves
\eqref{eq:finite-many-gap}.

For the infinite system, use the source's fixed homogeneous
empty-boundary exhaustion: retain a grouped interaction when all of its
declared support lies inside the site volume. This need not coincide with
retaining every individual boundary plaquette of a conventional finite
box, but it exhausts the homogeneous interaction. The external theorem
gives its canonical state $\omega_\infty$ and GNS triple
$(\mathcal H_\infty,\pi,\Omega)$, with~\eqref{eq:infinite-gap}.
Changing padded finite models is not used to assert this thermodynamic
existence.

It remains to justify the physical cyclic restriction. For a finite
vertex set $F$, let $\mathcal E_F$ be Haar averaging on local operator
algebras over the gauge transformations at $F$. It is a normal conditional
expectation in the underlying finite tensor product. Gauge-invariant
finite-volume vacua satisfy
\[
 \omega_\Lambda(B^*\mathcal E_F(A))
 =\omega_\Lambda(\mathcal E_F(B)^*A),\qquad
 \|\mathcal E_F(A)\Omega_\Lambda\|\le\|A\Omega_\Lambda\|.
\]
For fixed local $A,B,F$ these identities pass to the thermodynamic state.
Consequently
\[
 P_F\pi(A)\Omega=\pi(\mathcal E_F(A))\Omega
\]
defines an orthogonal projection on the GNS space. For increasing finite
sets $F$ exhausting all vertices the projections decrease and converge
strongly to a projection $P$. Averaging any local operator over all
vertices incident to its links already makes it globally gauge invariant.
It follows, by applying the projections to the dense set of local GNS
vectors, that
\begin{equation}
 \operatorname{Ran}P=
 \overline{\{\pi(A)\Omega:A\text{ local and gauge invariant}\}}.
 \label{eq:physical-gns}
\end{equation}
Finite-volume resolvents commute with gauge averaging. The weak resolvent
matrix-element convergence in Yarotsky's Theorem~3 passes this identity
to $P_F$ and the limiting resolvent, and then to $P$. Thus
\eqref{eq:physical-gns} reduces $H_\infty$, contains $\Omega$, and
inherits its simple vacuum and its gap. This argument does not require
assuming norm continuity of the gauge action on every local bounded
operator.
\end{proof}

The constant~\eqref{eq:threshold} is an existence threshold. No numerical
values for $c_1,c_2$ have been extracted, so the theorem does not establish
$X_d\le2$, or any other proposed numerical onset. It is also a
fixed-spacing, fixed-$d$ assertion. A limit with fixed 't~Hooft coupling
$Nx$ and growing $N$ eventually leaves the assumed window. Periodic tori
are not included by the open-subgraph argument. The simple vacuum in the
canonical GNS representation does not assert uniqueness among all
infinite-volume ground-state representations. No continuum scaling,
confinement observable, or relation to the scalar theory in
Section~\ref{sec:source-bridge} is supplied by this theorem.

\section{Benchmarks, reproducibility, and certification limits}
\label{sec:reproduce}

\subsection{Numerical checks}

The radial implementation generates equilibria rather than importing
stellar evolution profiles. For the $n=3$ models it uses
$G=\rho_c=\text{Lane--Emden length scale}=1$, $P_c=\pi$, and
$c^2=\pi/s$ when $s=P_c/(\rho_cc^2)>0$. It integrates the TOV or
Lane--Emden equations adaptively and assembles continuous piecewise-linear
Rayleigh--Ritz forms on a uniform radial mesh with eight-point Gaussian
quadrature per cell. Endpoints are not quadrature nodes, and no artificial
zero-displacement condition is imposed at the surface. Calculations use
binary64 arithmetic, not interval enclosures.

The generated Newtonian $n=3$ solution has radius approximately
$6.89684861937$ and Lane--Emden mass moment $2.01823595097$. At
$\Gamma=4/3$ it reproduces the homologous zero mode. At $\Gamma=1.5$,
the 256-cell value in Table~\ref{tab:benchmarks} is below the homologous
quotient; the latter exceeds it by approximately $20.43\%$. Nested
32-, 64-, 128-, and 256-cell trial spaces give decreasing computed Ritz
values. In exact arithmetic, nested Rayleigh--Ritz spaces provide upper
bounds and monotonicity; the floating calculation checks their numerical
realization and does not produce a certified lower bound.

\begin{table}[htbp]
\centering\small
\begin{tabular}{@{}p{0.45\linewidth}rr@{}}
\toprule
Benchmark & Comparison value & Computed value\\
\midrule
$n=3$, Newtonian, $\Gamma=1.5$
 & $Q_N[1]=0.5129469320$ & Ritz $0.4259235109$\\[3pt]
$n=3$, TOV, $s=0.001$, selected $\Gamma$
 & Homology $+2.41886\!\times\!10^{-7}$
 & Ritz $-2.42086\!\times\!10^{-7}$\\[3pt]
Single $\SU(2)$ plaquette, $x=2$, $b_N=2$
 & Lower bound $2$ & Ritz gap $3.1138638115$\\[3pt]
Single $\SU(3)$ plaquette, $x=2$, $b_N=2$
 & Lower bound $13/3$ & Ritz gap $5.1762835133$\\
\bottomrule
\end{tabular}
\caption{Selected checks, in the code's stated units. The TOV row uses
$\Gamma=1.335967888598$ and compactness $0.002325278673$. Radial Ritz
values use 256 cells. Character-matrix gaps are differences of finite
Ritz levels and are not certified lower bounds on the full gap. The
analytic bounds are proved independently of these computations.}
\label{tab:benchmarks}
\end{table}

The TOV row is a numerical illustration of a positive homologous value
coexisting with negative work on another admissible trial space.
Its small sign difference has not been interval-certified. The analytic
comparison direction and the residual obstruction
\eqref{eq:homology-residual} do not rely on that numerical sign.
Independent assembly of~\eqref{eq:gr-rayleigh} and
\eqref{eq:stable-gr-form} agrees in the tested model to about
$2.36\times10^{-12}$ in the largest matrix-entry difference.

For I14 the numerical check constructs the original rectangular
first-difference operator $B u=(u_{n+1}-3u_n/2)_n$, then forms
$B^TB$ and the scalar shifts. Thirty matrices, using
$L\in\{2,3,4,8,31,128\}$ and five parameter triples, agree with
\eqref{eq:sine-spectrum} to at most $6.22\times10^{-15}$. Exact symbolic
checks separately verify the source-action derivatives and flux
obstruction. These tests check constants and indexing; completeness and
sharpness follow from the proofs in Section~\ref{sec:discrete}.

The I15 script performs exact rational Casimir checks over a declared
finite family and constructs $\SU(2)$ and $\SU(3)$ character truncations
at multiple cutoffs. Its 159 recorded assertions cover normalization,
the displayed inequalities, cutoff comparisons, and failure of bare-vacuum
subtraction. Finite diagonalizations are not used to infer a universal
Casimir minimum or an infinite-space lower gap bound.

\subsection{Formal coverage}

The three Lean sources compile in the recorded project environment.
Their printed theorem dependencies contain only \texttt{propext},
\texttt{Classical.choice}, and \texttt{Quot.sound}; no custom axioms or
proof placeholders were introduced in these certificates. This describes
the checked declarations, not a formal verification of the complete
manuscript.

For I13, the existing scalar proof bodies were retained and their
interpretation corrected. The radial differential operator, domain
analysis, and post-Newtonian expansion remain analytic arguments in
prose. For I14, seventeen added declarations cover sharp uniform
inequalities, confinement, finite positivity, sine recurrence and energy
identities, and nonzero norms. The completeness argument and exact
finite spectral minimum are supplied in Section~\ref{sec:discrete}, not
claimed as fully formalized declarations. For I15, six added declarations
cover scalar comparisons, the Casimir floor expression, local perturbation
normalization, and energy rescaling. Peter--Weyl theory, min--max, the
external stability theorem, and the gauge/GNS transfer are not claimed
as Lean-formalized.

\subsection{Archive layout and reproduction}

Repository-relative source paths are used throughout the archive. The
consolidated directory is
\begin{center}\small
\path{real_research/reviews/spectral_spine_closure_2026_09_22/}.
\end{center}
Its \path{REPORT.md} records the combined scope; the principal derivations
are \path{i13/REPORT.md}, \path{i14/DERIVATION.md},
\path{i14/SOURCE_BRIDGE.md}, and \path{i15/PROOF.md}. The audit baseline
is commit \texttt{e3af62a453ca6625db8428f6d32957b27a7f1331}.
This baseline identifies the starting inputs, not a claim that it already
contains the later results. Final file hashes and executed verification
records identify the checked artifacts. The publication archive should be
used with its matching source files rather than an unpinned future branch.

The numerical and symbolic entry points are
\path{i13/radial_bridge.py}, \path{i13/symbolic_identities.py},
\path{i14/verify.py}, and \path{i15/verify.py}, relative to that
directory. Their adjacent contracts, results, manifests, and preserved
standard output specify commands, dependencies, parameter ranges,
resource limits, input hashes, and exit status. For example, from the
repository root the radial checks can be repeated as
\begin{quote}\small\ttfamily
python3 real\_research/reviews/spectral\_spine\_closure\_2026\_09\_22/\\
\hspace*{1em}i13/radial\_bridge.py --output /tmp/i13\_results.json
\end{quote}
where the displayed line break is typographical and the script path is a
single shell argument. The symbolic script accepts its output JSON path
as a positional argument. For all checks, the stored manifest is the
authoritative command record; paths can be redirected to a local output
directory without changing the mathematical inputs.

The Lean files are in \path{fable_independent_2026/lean_2026/}:
\begin{quote}\small
\path{I13_bhstar_pulsation.lean}\\
\path{I14_phantom_vacuum_wall.lean}\\
\path{I15_suN_lattice_gap.lean}
\end{quote}
From that directory each is checked with
\texttt{lake env lean} followed by its filename. The archive records
successful runs, with 17 numerical assertions and four symbolic identities
for I13, the I14 matrix and symbolic checks, and the I15 assertion suite.
The verification directory preserves combined compiler and manifest
summaries. Successful execution and matching hashes establish provenance
and the recorded checks; they do not independently prove every analytic
step.

The source records pin Saio et al. to version \path{2406.18040v1} and
Yarotsky to \path{math-ph/0411042v1}. They specify theorem and equation
locations and the application boundaries. Those
records do not claim that complete source PDFs were retained or hashed
locally. The mathematical claims use the pinned source versions listed
in the bibliography. Public repository links are provided in
References~\cite{archive,leanrepo,sources}.

\section{Consequences and remaining mathematical tasks}

The radial calculation supplies the missing variational comparison and
fixes the inertia and relativistic coefficient conventions. A stellar
mass cutoff still requires an actual equilibrium sequence, appropriate
equation of state, relevant inner and outer boundaries, and either the
exact relativistic test or a bounded post-Newtonian remainder. The
generated polytropes are useful benchmarks, not MESA or GENEC profile
data and not a black-hole-envelope mass sequence.

The discrete calculation closes the spectrum of the specified
constant-confinement model, including the distinction between a finite
box and a bound uniform in box size. The action calculation then identifies
why that theorem cannot presently select a physical acceleration
coupling: the cited source gives a different radial weight, no independently
derived negative $\kappa^2$ potential, and a logarithmic profile that is
only an approximate background for its full interpolant. Replacing an
algebraic coefficient does not resolve those operator-identification
requirements.

The lattice results provide two actual gap comparisons, one elementary
and one using an established external theorem. Both refer to the shifted
interacting ground state. The open-volume theorem is uniform in gauge
rank because the local Casimir normalization controls the interaction
norm independently of $N$. Evaluating a useful explicit $X_d$, reaching
$x=2$ uniformly in arbitrary open volume, and analyzing any continuum
limit are separate tasks. None is established by the positivity of the
single-plaquette scalar expression.

Together the examples show the practical role of a certificate's
interface: one must say which quadratic form is represented, in which
norm and domain, and whether the certified expression is an upper trial
value, a lower bound, or an exact eigenvalue. Those distinctions determine
which physical conclusions the mathematics can support.

\section*{License}

This manuscript is licensed under the Creative Commons Attribution 4.0
International license (\href{https://creativecommons.org/licenses/by/4.0/}{CC BY 4.0}).
The accompanying software supplement retains its stated GNU Affero General
Public License, version 3 or later (AGPL-3.0-or-later); the manuscript
license does not relicense that software.

\section*{AI assistance and review status}

OpenAI Codex assisted with drafting, symbolic and numerical code, source
inspection, and mathematical consistency checks for this manuscript and
its accompanying artifacts. Separate AI-assisted checks were used for
some operator and source mappings. Such checks are not independent human
peer review, and the disclosure does not imply that every computation or
derivation has been manually reviewed by the author. The executable
records and stated limits identify the evidence being presented. No AI
system is listed as an author.

\begin{thebibliography}{9}

\bibitem{saio}
H.~Saio, D.~Nandal, S.~Ekstr\"om, and G.~Meynet,
\emph{Linear adiabatic analysis for general relativistic instability in
primordial accreting supermassive stars},
arXiv:2406.18040v1 (26 June 2024).
\url{https://arxiv.org/abs/2406.18040v1}.
Equations (2)--(5) specify the radial source operator; equation (11)
provides the quoted $n=3$ benchmark.

\bibitem{yarotsky}
D.~A.~Yarotsky,
\emph{Quasi-particles in weak perturbations of non-interacting quantum
lattice systems}, arXiv:math-ph/0411042v1 (11 November 2004).
\url{https://arxiv.org/abs/math-ph/0411042v1}.
Theorems 1--3 and equations (1)--(6), pp.~2--4; proof review in Section~2.

\bibitem{ligterink}
N.~E.~Ligterink, N.~R.~Walet, and R.~F.~Bishop,
\emph{Towards a Many-Body Treatment of Hamiltonian Lattice SU(N) Gauge
Theory}, arXiv:hep-lat/0001028v1 (25 January 2000);
\emph{Annals of Physics} \textbf{284} (2000), 215--262.
\url{https://arxiv.org/abs/hep-lat/0001028v1}.
\href{https://doi.org/10.1006/aphy.2000.6070}{doi:10.1006/aphy.2000.6070}.

\bibitem{archive}
C.~P.~Zimmerman,
\emph{Spectral spine closure: derivations, computational contracts,
source records, and verification outputs}, accompanying research archive
(22 September 2026).
Repository directory
\path{real_research/reviews/spectral_spine_closure_2026_09_22/}.
\href{https://github.com/carlzimmerman/zimmerman-formula/tree/main/real_research/reviews/spectral_spine_closure_2026_09_22}{Public repository: spectral derivations and verification}.
Use the matching publication archive and its file hashes for reproduction.

\bibitem{leanrepo}
C.~P.~Zimmerman,
\emph{I13, I14, and I15 Lean certificates},
\path{fable_independent_2026/lean_2026/}, accompanying repository sources.
\href{https://github.com/carlzimmerman/zimmerman-formula/tree/main/fable_independent_2026/lean_2026}{Public repository: Lean source directory}.
Compilation records and final source hashes appear in Reference~\cite{archive}.

\bibitem{sources}
C.~P.~Zimmerman,
\emph{Project source actions and prior spectral calculations},
\path{deepseek_push/G155_sourced_eq.py},
\path{deepseek_push/G081_equilibrium_stability.py},
\path{deepseek_push/WAVEBOARD.md}, and
\path{deepseek_push/THE_THEORY.md}.
\href{https://github.com/carlzimmerman/zimmerman-formula/tree/main/deepseek_push}{Public repository: inspected project sources}.
These are provenance sources for the proposed physical identification;
the limitations established in Section~\ref{sec:source-bridge} apply.

\end{thebibliography}
\end{document}
